\subsection{交换子与对偶}

设 A 与 B 是 k-代数。回忆（III, §4, No.~3, 第 466 页）代数 A 与 B 上的双模是这样的 (A,B)-双模 P：由 A-模结构与 B-模结构导出的两个 k-模结构在它上面相合。为避免歧义，我们说 P 是 \((\mathrm{A},\mathrm{B})_{k}\)-双模。设 P 是一个 \((\mathrm{A},\mathrm{B})_{k}\)-双模。我们用 \(P^{*}\) 表示 P 的底层左 A-模的对偶 \(\operatorname{Hom}_{\mathrm{A}}(\mathrm{P},\mathrm{A})\)。它是一个 \((\mathrm{B},\mathrm{A})_{k}\)-双模（II, §1, No.~14, 第 226 页）；对 \(a \in A\)、\(b \in B\)、\(x \in P\) 及 \(x^{*} \in P^{*}\)，我们有

\[
\langle x, b x ^ {*} a \rangle = \langle x b, x ^ {*} \rangle a.\tag{1}
\]

我们用 \(_{s}A_{d}\) 表示视为 \((\mathrm{A},\mathrm{A})_{k}\)-双模（同前所引）的代数 A，并用 \(\Lambda:P\otimes_{B}P^{*}\to_{s}A_{d}\) 表示由下式确定的 \((\mathrm{A},\mathrm{A})_{k}\)-双模同态

\[
\Lambda (x \otimes x ^ {*}) = \langle x, x ^ {*} \rangle\tag{2}
\]

其中 \(x \in \mathrm{P}\) 且 \(x^{*} \in \mathrm{P}^{*}\) 。我们用 \(\widetilde{\mathrm{P}}\) 表示对偶 \(\mathrm{Hom}_{\mathrm{B}}(\mathrm{P},\mathrm{B})\)（\(\mathrm{P}\) 的底层右 B-模的）；它是一个 \((\mathrm{B},\mathrm{A})_{k}\)-双模。我们用 \(\widetilde{\Lambda}: \widetilde{\mathrm{P}} \otimes_{\mathrm{A}} \mathrm{P} \to {}_{s}\mathrm{B}_{d}\) 表示由下式确定的 \((\mathrm{B},\mathrm{B})_{k}\)-双模同态

\[
\widetilde {\Lambda} (\widetilde {x} \otimes x) = \langle \widetilde {x}, x \rangle\tag{3}
\]

其中 \(x \in P\)，\(\widetilde{x} \in \widetilde{P}\)。

现在，设映射 \(b \mapsto b_{\mathrm{P}}\) 是从 B 到 \(\operatorname{End}_{\mathrm{A}}(\mathrm{P})\) 的双射；这时它是从 B 到 \(\operatorname{End}_{\mathrm{A}}(\mathrm{P})\) 的反代数的同构。\(\mathbf{Z}\)-模的一个典范同态（从 \(\mathrm{P}^* \otimes_{\mathrm{A}} \mathrm{P}\) 到 \(\operatorname{End}_{\mathrm{A}}(\mathrm{P})\)；II, §4, No.~2, 第 271 页）于是确定了一个 \(\mathbf{Z}\)-模同态 \(\Theta: \mathrm{P}^* \otimes_{\mathrm{A}} \mathrm{P} \to \mathrm{B}\)，它由下式定义

\[
x \Theta (x ^ {*} \otimes y) = \langle x, x ^ {*} \rangle y\tag{4}
\]

其中 \(x, y \in \mathrm{P}\)，\(x^{*} \in \mathrm{P}^{*}\)。由 (1) 式，这个同态是 (B,B)-线性的，且我们有

\[
\Theta (x ^ {*} \otimes y) y ^ {*} = x ^ {*} \langle y, y ^ {*} \rangle\tag{5}
\]

其中 \(y \in \mathrm{P}\)，\(x^{*}, y^{*} \in \mathrm{P}^{*}\)。由 (4) 与 (5)，我们在 \((\mathrm{B}, \mathrm{B})_{k}\)-双模 \(\mathrm{P}^{*} \otimes_{\mathrm{A}} \mathrm{P}\) 中推得下列等式：

\[
(y ^ {*} \otimes x) \Theta (x ^ {*} \otimes y) = y ^ {*} \otimes \langle x, x ^ {*} \rangle y = y ^ {*} \langle x, x ^ {*} \rangle \otimes y = \Theta (y ^ {*} \otimes x) (x ^ {*} \otimes y)
\]

其中 \(x, y \in \mathrm{P}\)，\(x^{*}, y^{*} \in \mathrm{P}^{*}\)，从而

\[
s \Theta (t) = \Theta (s) t\tag{6}
\]

其中 \(s, t \in \mathrm{P}^* \otimes_{\mathrm{A}} \mathrm{P}\)。同样，对 \(x, y\)（在 \(\mathrm{P}\) 中）及 \(x^*, y^*\)（在 \(\mathrm{P}^*\) 中），我们由 (4) 与 (5) 在 \((\mathrm{A}, \mathrm{A})_k\)-双模 \(\mathrm{P} \otimes_{\mathrm{B}} \mathrm{P}^*\) 中推得下列等式：

\[
(x \otimes x ^ {*}) \langle y, y ^ {*} \rangle = x \otimes \Theta (x ^ {*} \otimes y) y ^ {*} = x \Theta (x ^ {*} \otimes y) \otimes y ^ {*} = \langle x, x ^ {*} \rangle (y \otimes y ^ {*}),
\]

从而

\[
u \Lambda (v) = \Lambda (u) v\tag{7}
\]

其中 \(u, v \in \mathrm{P} \otimes_{\mathrm{B}} \mathrm{P}^{*}\)。

对每个元素 \(x^{*}\)（属于 \(P^{*}\)），把从 P 到 B 的 B-线性映射 \(x \mapsto \Theta(x^{*} \otimes x)\) 记作 \(\sigma(x^{*})\)。这样我们就定义了一个映射 \(\sigma\)（从 \(P^{*}\) 到 \(\widetilde{P}\)），它是 (B, A)-线性的，且按定义满足

\[
\Theta (x ^ {*} \otimes y) = \langle \sigma (x ^ {*}), y \rangle\tag{8}
\]

其中 \(x^{*} \in P^{*}\)，\(y \in P\)。由 \(\widetilde{\Lambda}\) 的定义，我们因此有

\[
\Theta = \widetilde {\Lambda} \circ (\sigma \otimes 1 _ {\mathrm{P}}).\tag{9}
\]

设映射 \(a \mapsto a_{\mathrm{P}}\)（从 A 到 \(\operatorname{End}_{\mathrm{B}}(\mathrm{P})\)）是双射；这时它是一个代数同构。类似地，我们用下式定义一个 \((\mathrm{A}, \mathrm{A})_k\)-双模的同态 \(\widetilde{\Theta}: \mathrm{P} \otimes_{\mathrm{B}} \widetilde{\mathrm{P}} \to {}_s \mathrm{A}_d\)：

\[
\widetilde {\Theta} (x \otimes \widetilde {y}) y = x \langle \widetilde {y}, y \rangle\tag{10}
\]

其中 \(x, y \in \mathrm{P}\)，\(\widetilde{y} \in \widetilde{\mathrm{P}}\)。我们还用下式定义一个 \((\mathrm{B}, \mathrm{A})_k\)-双模的同态 \(\widetilde{\sigma}: \widetilde{\mathrm{P}} \to \mathrm{P}^*\)：

\[
\widetilde {\Theta} (x \otimes \widetilde {y}) = \langle x, \widetilde {\sigma} (\widetilde {y}) \rangle\tag{11}
\]

其中 \(x\in \mathrm{P},\widetilde{y}\in \widetilde{\mathrm{P}}\)。我们有

\[
\widetilde {\Theta} = \Lambda \circ (1 _ {P} \otimes \widetilde {\sigma}).\tag{12}
\]

命题 1. —— 设映射 \(b \mapsto b_{P}\)（从 B 到 \(\operatorname{End}_{\mathrm{A}}(\mathrm{P})\)）与映射 \(a \mapsto a_{P}\)（从 A 到 \(\operatorname{End}_{\mathrm{B}}(\mathrm{P})\)）都是双射。则 \(\sigma\) 与 \(\widetilde{\sigma}\) 是互逆的同构，且我们有关系式

\[
\Lambda = \widetilde {\Theta} \circ (1 _ {\mathrm{P}} \otimes \sigma),\tag{13}
\]

\[
\widetilde {\Lambda} = \Theta \circ (\widetilde {\sigma} \otimes 1 _ {\mathrm{P}}).\tag{14}
\]

对 \(x \in \mathrm{P}\)、\(x^{*} \in \mathrm{P}^{*}\) 及 \(y \in \mathrm{P}\)，由关系式 (4)、(8)、(10) 与 (11)，我们有

\[
\langle x, x ^ {*} \rangle y = x \Theta (x ^ {*} \otimes y) = x \langle \sigma (x ^ {*}), y \rangle = \widetilde {\Theta} (x \otimes \sigma (x ^ {*})) y = \langle x, \widetilde {\sigma} (\sigma (x ^ {*})) \rangle y.\tag{15}
\]

同样，对 \(x \in P\)、\(\widetilde{y} \in \widetilde{P}\) 及 \(y \in P\)，我们有

\[
x \langle \widetilde {y}, y \rangle = \widetilde {\Theta} (x \otimes \widetilde {y}) y = \langle x, \widetilde {\sigma} (\widetilde {y}) \rangle y = x \Theta (\widetilde {\sigma} (\widetilde {y}) \otimes y) = x \langle \sigma (\widetilde {\sigma} (\widetilde {y})), y \rangle .\tag{16}
\]

由这些假设，P 作为 A-模与作为 B-模都是忠实的。由关系式 (15) 与 (16)，我们分别推得 \(\widetilde{\sigma} \circ \sigma = 1_{P^{*}}\) 与 \(\sigma \circ \widetilde{\sigma} = 1_{\widetilde{P}}\)。关系式 (13) 与 (14) 于是分别由 (12) 与 (9) 得出。

注. —— 1) 设映射 \(b \mapsto b_{\mathrm{P}}\)（从 B 到 \(\operatorname{End}_{\mathrm{A}}(\mathrm{P})\)）是双射。则

\begin{enumerate}
\def\labelenumi{\alph{enumi})}
\item
  B-模 P 可与 P 的反模等同；因此它是忠实的且平衡的；
\item
  映射 \(a \mapsto a_{P}\)（从 A 到 \(\operatorname{End}_{\mathrm{B}}(\mathrm{P})\)）是双射，当且仅当 A-模 P 是忠实的且平衡的。
\end{enumerate}

\begin{enumerate}
\def\labelenumi{\arabic{enumi})}
\setcounter{enumi}{1}
\tightlist
\item
  在命题 1 的假设下，A-模 P 是平衡的；由于环 \(A_{P}\) 与 \(A_{P}^{\prime}\) 有相同的中心（VIII, 第 77 页），存在一个同构 \(\varphi\)（从环 A 的中心 \(Z(A)\) 到环 B 的中心 \(Z(B)\)），它由关系式 \(\varphi(z)_{\mathrm{P}} = z_{\mathrm{P}}\)（对 \(z \in \mathrm{Z}(A)\)）确定。此外，\((\mathrm{A}, \mathrm{B})_{k}\)-双模 P 的自同态是诸位似 \(z_{P}\)，其中 z 取遍 \(Z(A)\)；\((\mathrm{A}, \mathrm{B})_{k}\)-双模 P 的自同构是诸位似 \(z_{P}\)，其中 z 在 \(Z(A)\) 中可逆。
\end{enumerate}
% semantic-fragments: legacy-input-seam
\relax{}
